\section{Method}
\label{sec:method}

Given a pre-trained text-to-image diffusion model and a set of target concepts to erase (\eg, nudity, violence), \ours{} intervenes at inference time through three stages: \textbf{When} to activate safety guidance, \textbf{Where} to localize the intervention spatially, and \textbf{How} to steer the denoising trajectory toward safe content. The whole pipeline is driven by a lightweight \emph{concept pack}---a family-grouped collection of exemplar prompts and images---so that adding or combining concepts never requires re-training. Figure~\ref{fig:overview} illustrates the full pipeline.

\subsection{Preliminaries and Concept Packs}
\label{subsec:prelim}

\paragraph{Classifier-free guidance.} In CFG~\citep{ho2022classifier}, the noise prediction at step $t$ is
\begin{equation}
    \hat{\epsilon}_t = \epsilon_\varnothing + s \cdot (\epsilon_\theta(z_t, c) - \epsilon_\varnothing),
\end{equation}
where $\epsilon_\varnothing = \epsilon_\theta(z_t, \varnothing)$ is the unconditional prediction, $\epsilon_\theta(z_t, c)$ is conditioned on prompt $c$, and $s$ is the guidance scale. We call $d_c(t) = \epsilon_\theta(z_t, c) - \epsilon_\varnothing$ the \emph{concept direction} at step $t$.

\paragraph{Family-grouped concept pack.}
Unsafe concepts are rarely monolithic: \emph{nudity} spans \textit{naked body}, \textit{sexual act}, \textit{exposed body parts}, and \textit{suggestive pose}; \emph{violence} spans \textit{blood/gore}, \textit{weapon attack}, \textit{brutal action}, and \textit{dismemberment}; \emph{hate} covers racial, religious, gender-based and political axes. Pooling all exemplars into one averaged vector collapses this diversity and yields a weak, generic direction that under-fits adversarial prompts.
We therefore adopt \emph{family-grouped exemplars} as the first-class representation of \ours{}. For every concept we define $F$ semantically disjoint \emph{families} (default $F{=}4$), and for each family $f\in\{1,\dots,F\}$ we pair $K$ target exemplars with $K$ anchor exemplars (default $K{=}4$). A concept pack therefore materializes to
\begin{equation}
    \mathcal{P} = \big\{(c^{\text{tgt}}_{f,k},\, c^{\text{anc}}_{f,k}, I^{\text{tgt}}_{f,k}, I^{\text{anc}}_{f,k})\big\}_{f=1,\dots,F;\ k=1,\dots,K} \cup \mathcal{W}_f,
    \label{eq:pack}
\end{equation}
where $\mathcal{W}_f$ are family-level target keywords (\eg, ``naked'', ``exposed'' for \emph{naked body}). All text- and image-side quantities used below are indexed by family.

\paragraph{Family-level embeddings.}
Given $\mathcal{P}$, we precompute \emph{once}: (i) family-mean CLIP embeddings $\bar{e}_f = \tfrac{1}{K}\sum_k \mathrm{CLIP}(I^{\text{tgt}}_{f,k})$ and their anchor counterparts $\bar{a}_f$; (ii) family-level target and anchor text directions $d^{\text{tgt}}_f(t) = \tfrac{1}{K}\sum_k \big[\epsilon_\theta(z_t, c^{\text{tgt}}_{f,k}) - \epsilon_\varnothing\big]$ and $d^{\text{anc}}_f(t)$. No gradients or parameter updates are needed; adding a new concept amounts to adding one directory with $4F$ prompts and $4F$ images.

\subsection{When: Concept Alignment Score (CAS)}
\label{subsec:when}

Not every denoising step needs safety intervention. Early steps establish global composition while later steps refine details; unsafe content may emerge only at specific timesteps. Applying guidance indiscriminately wastes compute and risks degrading benign content.

We gate intervention by a family-aware \textbf{Concept Alignment Score (CAS)}:
\begin{equation}
    \text{CAS}(t) \;=\; \max_{f=1,\dots,F} \cos\!\big(d_\text{prompt}(t),\; d^{\text{tgt}}_f(t)\big),
    \qquad d_\text{prompt}(t) = \epsilon_\theta(z_t, c) - \epsilon_\varnothing,
    \label{eq:cas}
\end{equation}
so that the denoising direction is compared against \emph{every} sub-category independently; the highest match wins. Safety guidance fires only when $\text{CAS}(t) > \tau$ (typically $\tau{\in}[0.4, 0.6]$). Benign prompts---whose directions do not align with any family---proceed without intervention, which sharply reduces false positives.

\paragraph{Sticky CAS.} Once $\text{CAS}(t_0) > \tau$ at any step $t_0$, we latch activation for all $t > t_0$. This prevents flicker between active/inactive states, which empirically produces inconsistent renderings.

\subsection{Where: Dual-Probe Spatial Localization}
\label{subsec:where}

When CAS triggers intervention, the next question is \emph{where} in the spatial domain the unsafe concept is forming. \ours{} combines complementary \emph{text} and \emph{image} probes, each evaluated per family and then unioned.

\paragraph{Text probe.} Cross-attention maps in the U-Net bind each spatial location to specific text tokens. Using the family keyword set $\mathcal{W}_f$, we extract $A_w\in\mathbb{R}^{H\times W}$ for $w\in\mathcal{W}_f$ and form
\begin{equation}
    M^f_\text{text}(i,j) = \sigma\!\Big(\alpha \cdot \big(\textstyle\max_{w\in\mathcal{W}_f} A_w(i,j) - \theta_\text{text}\big)\Big),
    \label{eq:text_probe}
\end{equation}
with sigmoid sharpness $\alpha$ and attention threshold $\theta_\text{text}$.

\paragraph{Image probe.} Text probes fail when adversarial prompts avoid explicit keywords. A CLIP image probe bypasses the prompt: at each step we extract intermediate U-Net features $f(i,j)$ and score them against the family-mean embeddings $\{\bar{e}_f\}$,
\begin{equation}
    M^f_\text{img}(i,j) = \sigma\!\Big(\alpha \cdot \big(\cos(f(i,j), \bar{e}_f) - \theta_\text{img}\big)\Big).
    \label{eq:img_probe}
\end{equation}
This flags visually concept-relevant regions regardless of how the prompt is phrased.

\paragraph{Probe fusion.}
We first union over families (a location is flagged if it matches \emph{any} sub-category) and then union across modalities:
\begin{equation}
    M_\text{text} = \max_f M^f_\text{text},\qquad M_\text{img} = \max_f M^f_\text{img},\qquad
    M = \min\!\big(1,\, M_\text{text} + M_\text{img}\big).
    \label{eq:fusion}
\end{equation}
Union fusion ensures that regions detected by \emph{either} probe---or any family---are suppressed, providing robustness against both explicit and metaphorical adversarial prompts. Text-only and image-only modes are retained for ablation.

\subsection{How: Spatially Masked Safe Guidance}
\label{subsec:how}

Given the spatial mask $M$ and the family decomposition, we modify $\hat{\epsilon}_t$ \emph{only} inside $M$. We use a standard mask-gated CFG adjustment and treat the two variants below as simple engineering choices rather than headline contributions:
\begin{align}
    \text{Anchor inpainting:}&\quad \hat{\epsilon}_\text{safe} = \hat{\epsilon}_t - s_\text{safe} \cdot M \cdot \Big(\tfrac{1}{F}\textstyle\sum_{f} d^{\text{tgt}}_f(t)\Big), \label{eq:anchor_inpaint}\\
    \text{Hybrid guidance:}&\quad \hat{\epsilon}_\text{safe} = \hat{\epsilon}_t + s_t \cdot M \cdot \Big(\tfrac{1}{F}\textstyle\sum_{f} \big[d^{\text{anc}}_f(t) - d^{\text{tgt}}_f(t)\big]\Big). \label{eq:hybrid}
\end{align}
Hybrid pushes the trajectory from target toward family-paired anchors (\eg, \textit{exposed parts}$\to$\textit{clothed counterpart}, \textit{weapon attack}$\to$\textit{calm interaction}); anchor inpainting only subtracts the target direction. Both act only inside the mask, so the rest of the image is untouched.

\subsection{Training-Free, User-Customizable Exemplars}
\label{subsec:training_free}

The central practical advantage of \ours{} is that \emph{everything about a concept lives in the exemplar pack}---no weights are modified, no auxiliary classifier is trained, no dataset is collected. Concretely:

\textbf{(1) Training-free.} All quantities used by CAS (Eq.~\ref{eq:cas}), the probes (Eq.~\ref{eq:text_probe}--\ref{eq:img_probe}), and the guidance (Eq.~\ref{eq:anchor_inpaint}--\ref{eq:hybrid}) are either precomputed once from exemplars or computed on the fly from the frozen UNet and a frozen CLIP. In contrast, ESD/SDD/RECE/SDErasure require per-concept fine-tuning ($\sim$20--30 min each, Appendix~\ref{app:compute}) and thereby risk catastrophic forgetting on benign prompts.

\textbf{(2) User-customizable, few-shot adaptation.} Because an exemplar pack is just $F\cdot K$ prompts plus $F\cdot K$ images (16 of each in our default), the user fully controls what ``unsafe'' means for their deployment. Harder adversarial distributions can be addressed by \emph{few-shot adaptation of the pack}: if MMA-style metaphorical prompts~\citep{zhang2024metaphor} are the target, one can drop a handful of MMA-style exemplars into the pack and re-precompute $\bar{e}_f$ and $d^{\text{tgt}}_f$ in seconds---no gradient step required. Conversely, removing a family (or keyword) narrows the scope with the same ease. This makes the method an \emph{editable} safety layer rather than a fixed artifact.

\textbf{(3) Simple but effective.} A consequence of (1)--(2) is that our method is mechanically light: text probes alone work very well on explicitly-named concepts (nudity), image probes alone pick up concepts with distinctive visual signatures (violence, disturbing), and the union covers both. We found that this simple combination is hard to beat within the training-free regime; complexity is not required, adaptability is.

Section~\ref{subsec:ablation} validates each of these claims empirically.

\subsection{Multi-Concept Erasure}
\label{subsec:multi}

\ours{} supports simultaneous multi-concept erasure by construction. Given $C$ target concept packs, each with its own CAS threshold $\tau_c$, spatial mask $M_c$ (Eq.~\ref{eq:fusion}) and guidance term $g_c$:
\begin{equation}
    \hat{\epsilon}_\text{safe} = \hat{\epsilon}_t + \sum_{c=1}^{C} \mathbb{1}\!\left[\text{CAS}_c(t) > \tau_c\right] \cdot s_c \cdot M_c \cdot g_c(z_t).
    \label{eq:multi}
\end{equation}
Each concept contributes $F$ family-level directions and $F$ family-level masks that are independently localized, so concepts only interact where their masks \emph{and} family directions overlap---a much smaller region than with a single pooled direction per concept. Adding a new concept means adding one exemplar pack; no retraining, no model surgery.

\paragraph{Intersection-only suppression (preliminary).} When two concept masks overlap (\eg, nudity and hate in a human-subject region), naively summing their guidance can collapse the content in the shared region. A simple mitigation is to apply guidance only on the \emph{intersection} $\bigcap_c M_c$ for overlapping concepts, leaving concept-specific regions to their own packs. We treat this as a lightweight safeguard and report preliminary numbers; a more principled treatment of multi-concept interference is left for future work.
